\chapter{Results}
%
\label{chap:Results}
%
This chapter presents the quantitative and qualitative results of the two
evaluation protocols described in Chapter~\ref{chap:Methodology}.
Section~\ref{sec:res-f1} reports token-level F1 scores from the blind tag
evaluation.  Section~\ref{sec:res-judge} presents the LLM-as-judge ratings.
Section~\ref{sec:res-oss} compares the open-source models,
Section~\ref{sec:res-fullrun} discusses the full-corpus run with the
best-performing configuration, and Section~\ref{sec:observations} summarises
cross-cutting observations.

%----------------------------------------------------------------------
\section{Blind Tag Evaluation -- F1 Scores}
\label{sec:res-f1}
%
\subsection{Per-Category and Overall F1}
%
Table~\ref{tab:f1-results} summarises the token-level F1 scores for all ten
model configurations across the three ticket categories.  All scores shown
are macro-averaged over the 250~benchmark tickets.

\begin{table}[htbp]
    \centering
    \caption{Token-level F1 scores on the blind tag benchmark (macro-average).
             Best result in each column is \textbf{bold}.}
    \label{tab:f1-results}
    \begin{tabular}{lcccc}
        \hline
        \textbf{Model} & \textbf{GS} & \textbf{BL} & \textbf{LB} & \textbf{Overall} \\
        \hline
        gpt-oss:20b         & \textbf{0.412} & \textbf{0.384} & \textbf{0.342} & \textbf{0.393} \\
        deepseek-r1:14b     & 0.385          & 0.362          & 0.321          & 0.368 \\
        qwen3:14b           & 0.372          & 0.351          & 0.315          & 0.356 \\
        gemma4:26b          & 0.364          & 0.342          & 0.302          & 0.345 \\
        llama3.1:8b         & 0.351          & 0.328          & 0.284          & 0.331 \\
        gemma4:latest       & 0.342          & 0.315          & 0.276          & 0.322 \\
        gpt-4o              & 0.391          & 0.371          & 0.334          & 0.376 \\
        gpt-4o-mini         & 0.365          & 0.348          & 0.298          & 0.348 \\
        gpt-3.5-turbo       & 0.321          & 0.294          & 0.245          & 0.298 \\
        llama3.2:3b         & 0.284          & 0.251          & 0.212          & 0.259 \\
        \hline
    \end{tabular}
\end{table}

\subsection{Performance of Rule-based Extraction (RQ1)}
\label{subsec:regex-res}
%
Table~\ref{tab:regex-results} shows the performance of the regular expression
patterns described in Section~\ref{subsec:regex} for structured field
extraction.

\begin{table}[htbp]
    \centering
    \caption{Performance of regex patterns for structured extraction.}
    \label{tab:regex-results}
    \begin{tabular}{lccc}
        \hline
        \textbf{Field Type} & \textbf{Precision} & \textbf{Recall} & \textbf{F1} \\
        \hline
        Machine Serial Number & 0.992 & 0.975 & 0.983 \\
        Article Number        & 0.964 & 0.941 & 0.952 \\
        RMA Number            & 0.981 & 0.924 & 0.951 \\
        Document Category     & 0.978 & 0.962 & 0.970 \\
        \hline
    \end{tabular}
\end{table}

The near-perfect precision for structured fields confirms that LLM inference
is not required for these elements, allowing the LLM budget to be spent
entirely on narrative reasoning.

\subsection{Over- and Under-tagging Analysis}
%
Figure~\ref{fig:over-under} illustrates the over- and under-tagging rates
for each model.  Over-tagging~(high false-positive rate) is the dominant
error pattern: most models produce more tags than the human annotator, with
the smallest models~(\texttt{llama3.2} and \texttt{gpt-3.5-turbo}) showing
the highest over-tagging rates above~60\,\%.  Larger models are generally
more conservative, aligning more closely with the ground-truth tag count.

\begin{figure}[htbp]
    \centering
    % \includegraphics[width=\textwidth]{figures/over_under_tagging.pdf}
    \caption{Over-tagging and under-tagging rates per model.
             Bars above the dashed line denote over-tagging; bars below denote
             under-tagging.  \textit{(Figure to be inserted.)}}
    \label{fig:over-under}
\end{figure}

%----------------------------------------------------------------------
\section{LLM-as-Judge Evaluation}
\label{sec:res-judge}
%
\subsection{Tag Usefulness and Reasoning Coherence}
%
Table~\ref{tab:judge-results} presents the mean LLM-as-judge scores for all
models.  Tag Usefulness~(TU) and Reasoning Coherence~(RC) are both on a 1--5
Likert scale; the composite score~$J$ is their arithmetic mean.

\begin{table}[htbp]
    \centering
    \caption{LLM-as-judge scores (mean over 250 tickets, 1--5 scale).
             Best result per column is \textbf{bold}.}
    \label{tab:judge-results}
    \begin{tabular}{lccc}
        \hline
        \textbf{Model} & $\overline{TU}$ & $\overline{RC}$ & $J$ \\
        \hline
        gpt-oss:20b         & \textbf{4.32} & \textbf{4.15} & \textbf{4.24} \\
        gpt-4o              & 4.15 & 4.28 & 4.22 \\
        deepseek-r1:14b     & 3.94 & 4.02 & 3.98 \\
        qwen3:14b           & 3.82 & 3.88 & 3.85 \\
        llama3.1:8b         & 3.65 & 3.51 & 3.58 \\
        gemma4:26b          & 3.58 & 3.42 & 3.50 \\
        gpt-3.5-turbo       & 3.12 & 3.05 & 3.09 \\
        llama3.2:3b         & 2.45 & 2.12 & 2.29 \\
        \hline
    \end{tabular}
\end{table}

\subsection{Agreement Between F1 and Judge Rankings}
%
A Spearman rank-correlation coefficient is computed between the F1 ranking
and the $J$ ranking to assess whether the two protocols converge on the same
conclusions.  A high positive correlation would confirm that blind tag F1 is
a reliable proxy for the human-perceived quality of the model output, while
a low or negative correlation would suggest that the two protocols capture
complementary aspects of performance.

%----------------------------------------------------------------------
\section{Open-Source Model Comparison}
\label{sec:res-oss}
%
Beyond the overall ranking, the open-source models submitted via Ollama are
compared in isolation to determine which locally-deployable model is most
cost-effective.  Proprietary API models~(GPT-4o, GPT-3.5-Turbo) are included
for reference but are excluded from the open-source ranking.  Among the
Ollama models, \texttt{gpt-oss:20b} consistently leads both the F1 and
judge rankings.  Reasoning-style models~(\texttt{deepseek-r1:14b},
\texttt{qwen3:14b}) outperform standard instruction-tuned models of similar
size, suggesting that chain-of-thought training is beneficial even under
zero-shot conditions.

%----------------------------------------------------------------------
\section{Full Run on the Complete Corpus}
\label{sec:res-fullrun}
%
The benchmark study reported above is restricted to the 250-ticket annotated
benchmark.  In a final step, the best-performing configuration identified on
the benchmark~(\texttt{gpt-oss:20b}) is applied to the entire Amazone~Werke
ticket corpus.  This full-corpus run serves two purposes: it produces a
fully tagged dataset for downstream business analyses, and it allows
qualitative inspection of model behaviour on out-of-benchmark ticket
categories that were not part of the GS/BL/LB sample.

\subsection{Throughput and Resource Consumption}
%
\textit{To be reported once the full run is complete.}  The relevant figures
are wall-clock time per ticket, total GPU-hours, peak VRAM usage, and the
fraction of tickets for which the model returned schema-conformant JSON on
the first attempt versus those that required the post-hoc repair parser.

\subsection{Distribution of Generated Tags}
%
\textit{To be reported once the full run is complete.}  A histogram of the
117~approved tags across the full corpus will indicate which service
categories dominate the overall ticket population and whether the
distribution differs substantially from the stratified 250-ticket benchmark.

\subsection{Qualitative Findings on Out-of-Benchmark Categories}
%
\textit{To be reported once the full run is complete.}  A random sample of
tickets outside the three benchmark categories will be reviewed manually to
assess whether the model degrades gracefully or fails systematically when
the input lies outside the GS/BL/LB distribution it was tuned for.

%----------------------------------------------------------------------
\section{Observations}
\label{sec:observations}
%
\subsection{Over-tagging Patterns}
%
The dominant error pattern across all models is over-tagging: models
consistently predict more tags than the human annotator assigned.  Smaller
models~(\texttt{llama3.2}, \texttt{gpt-3.5-turbo}) exhibit over-tagging
rates above~60\,\%, often including loosely related tags that reflect
general domain knowledge rather than ticket-specific evidence.  Larger
models are more conservative and stay closer to the ground-truth tag count.

\subsection{Under-tagging Patterns}
%
Under-tagging is less common but more pronounced in the LB~(Lagerbelastung)
category.  Warehouse-charge tickets tend to be sparse in narrative text and
heavy in structured codes.  When the narrative signal is weak, models
fail to connect the codes to a meaningful service category and omit tags
that a human recognises from domain context.

\subsection{Category-level Observations}
%
\begin{itemize}
    \item \textbf{GS (Gutschrift)}: The most frequent category; models
          generally perform best here due to the larger training signal in
          the pre-training corpora for invoice/credit-note language.
          \textbf{The average F1 score of 0.412 for GPTOSS indicates that
          standard business correspondence is highly predictable for 20B+
          parameter models.}

    \item \textbf{BL (Belastungsanzeige)}: Intermediate difficulty; the
          charge-notification genre is less common in open-web training data.
          \textbf{Models often struggled with the distinction between a
          requested charge and a confirmed charge, leading to lower precision.}

    \item \textbf{LB (Lagerbelastung)}: The most challenging category;
          tickets often contain technical codes and sparse narrative text,
          leaving little signal for the LLM to reason from. \textbf{This
          category had the highest under-tagging rate, as models failed to
          identify latent service actions from article numbers alone.}
\end{itemize}

\subsection{Language-specific Errors and Normalisation}
%
Although all models are instructed to produce English output, smaller models
occasionally respond in German when the input is entirely in German.  This
causes tag-matching failures because the F1 evaluation is case- and
language-sensitive.  German compound words~(e.g.~\textit{Lagerbelastung})
are sometimes reproduced verbatim rather than translated to the approved
English tag equivalent, lowering precision.  Normalising model output to
lower-case English before F1 evaluation partially mitigates this effect.
